\section{Smaller presentations and a printed specialization}\label{sec:smaller}

A generator with poles distributed over the fiber at infinity gives a
much smaller coefficient height for an equation of the same bidegree.
We first explain the degree bound, then prove the change of generator by
exact resultant identities and give a compact integral specialization.
For a nonzero rational polynomial, its \emph{primitive integer height}
is the maximum number of decimal digits of the absolute values of its
coefficients after clearing denominators and dividing out their common
integer factor. All heights below use this convention.

\Needspace{15\baselineskip}
\subsection{Coefficient degrees from pole orders}
\begin{lemma}\label{lem:polebound}
Suppose $t:Z\to\PP^1$ has degree $n$, and $x$ generates $\Q(Z)/\Q(t)$
with all its poles above $\infty$. At each geometric point $P$ above
$\infty$, let $e_P$ be the ramification index and let $k_P$ be the pole
order of $x$, or $0$ if $x$ is regular there. Form a list
$\lambda_1\geq\cdots\geq\lambda_n$ by repeating $k_P/e_P$ exactly $e_P$
times for each $P$. Then the minimal polynomial
\[
 F(X,t)=X^n+\sum_{i=1}^{n}c_i(t)X^{n-i}
\]
has $c_i(t)\in\Q[t]$ and
\[
 \deg_t c_i\leq\left\lfloor\lambda_1+\cdots+\lambda_i\right\rfloor.
\]
\end{lemma}
\begin{proof}
The function $x$ is regular above the affine $t$-line, so is integral over
$\Q[t]$. Its minimal polynomial therefore has coefficients in $\Q[t]$.
At infinity, each of the $e_P$ local sheets associated with $P$ grows at most like
$t^{k_P/e_P}$. Every term in the $i$th elementary symmetric function is
therefore bounded by a constant times $|t|^{\lambda_1+\cdots+\lambda_i}$.
Since $c_i$ is a polynomial, its degree satisfies the asserted bound.
\end{proof}

For a generator with pole divisor $3P_0+P_4$ as in
\eqref{eq:polefiber}, the list has thirteen entries $3/13$ and fifty-two entries
$1/13$, so
\[
 \deg_t c_i\leq
 \left\lfloor\frac{3\min(i,13)+\max(i-13,0)}{13}\right\rfloor\leq7.
\]
Thus such a generator gives an equation of $t$-degree at most $7$.
The exact coefficients of $F_3$ attain this bound. The affine substitution
$U=3(t+1)$ preserves the degree.

\subsection{The smaller family and its exact bridge}\label{app:smaller}

The polynomial $F_3(Y,U)$ has $318$ nonzero terms and primitive integer
height $47$ digits, compared with $249$ digits for $f$, although $F_3$ is
the denser polynomial. In the coordinate
$U=3(t+1)$ its leading terms are
\begin{equation}\label{eq:F3-leading}
\begin{split}
F_3(Y,U)={}&Y^{65}+78Y^{64}+2652Y^{63}+48880Y^{62}+457290Y^{61}\\
 &+(296636-13312U)Y^{60}+\cdots.
\end{split}
\end{equation}
The full file \texttt{data/F3\_L3P0P4.json} uses the exponent-pair encoding
in \cref{app:data}. Its generator comes from the two-dimensional space
$L(3P_0+P_4)$ in the discovery model. The following criterion proves its
relation to $f$.

\Needspace{11\baselineskip}
\begin{lemma}[Exact change of generator]\label{lem:bridge}
Let $K$ be a field of characteristic zero, and let $F(X),P(Y)$ be monic
irreducible polynomials of the same degree in $K[X]$ and $K[Y]$.
For $\Psi\in K[X,Y]$, put
\[
 A(X)=\Res_Y(P(Y),\Psi(X,Y)),\qquad
 B(Y)=\Res_X(F(X),\Psi(X,Y)).
\]
If $F\mid A$, $F^2\nmid A$, $P\mid B$, and $P^2\nmid B$, then
$\Psi(\alpha,\beta)=0$ gives a Galois-equivariant bijection between the
roots of $F$ and $P$. Corresponding roots generate the same extension of
$K$, and the two splitting fields coincide.
\end{lemma}
\begin{proof}
Since $P$ is monic, $A(X)=\prod_{P(\beta)=0}\Psi(X,\beta)$.
Irreducibility and separability of $F$, together with its multiplicity one
in $A$, show that each root $\alpha$ is a simple root of $A$. Exactly one
factor of the product therefore vanishes at $\alpha$. The same argument
with $B$ gives uniqueness in the reverse direction. The relation is
 defined over $K$, so the bijection commutes with every automorphism of an
algebraic closure of $K$. Corresponding roots have equal stabilizers,
proving the field assertions.
\end{proof}

Apply the lemma over $K=\Q(t)$ with $F=f(X,t)$ and
$P=F_3(Y,3(t+1))$. The supplied bridge file records a relation $\Psi$ of
bidegree $(7,7)$ and the affine change of generator used to express it.
The verifier checks the relation with that change applied; the public
$F_3$ agrees after the parameter substitution $U=3(t+1)$.
Irreducibility of $P$ is checked at $t=0$. The two resultant remainders
have $t$-degree at most $294$ and $154$: division by a monic polynomial
preserves the weights supplied by the pole bounds. Their exact vanishing
at $295$ and $155$ rational values, respectively, proves the two
divisibilities. Coprimality of each quotient with its divisor at $t=0$
proves multiplicity one.

It follows that $F_3$ has the same splitting field as $f$ over $\Q(t)$.
Writing $d(s)=3(s^3-7s^2+12s-5)$, its polynomial pullback is
\begin{equation}\label{eq:G3}
 G_3(Y,s)=d(s)^7F_3(Y,3(T(s)+1)).
\end{equation}
It has bidegree $(65,21)$. The stored pullback uses $s\mapsto s+1$ and has
primitive integer height $52$ digits; its exact substitution identity is
checked separately.
This proves the generic assertions of \cref{cor:smaller}.

\subsection{A printed specialization at \texorpdfstring{$s=1$}{s=1}}\label{sec:s1}

For $s=1$, so $t=1/3$ and $U=4$, the polynomial $F_3(Y,4)$ has height
$48$ digits. Its separable reductions have types $13^5$ at $11$ and
$1^2 7^9$ at $17$. This is a separable specialization of the regular
$N$-cover defined by $F_3(Y,3(T(s)+1))$. Its group embeds in $N$, and the
maximal-subgroup argument of \cref{sec:specialization} gives Galois group
$\Sz(8)$.

The monic polynomial $P_1$ in
\texttt{data/P\_s1\_number\_field.json}, printed in \cref{app:polynomial},
has height $43$ digits, compared with $68$ digits for the reduced generator
at $s=0$. The file supplies $A_1\in\Q[X]$ with
$F_3(A_1(X),4)\equiv0\pmod{P_1(X)}$. Exact irreducibility and this identity
show that $P_1$ defines the same specialization field. Certified conjugate
root matching gives its signature $(1,32)$, completing
\cref{cor:smaller}. The program \texttt{scripts/check\_review\_evidence.py}
checks these statements and compares all printed coefficients with the file.
